\documentclass[11pt]{article}
\usepackage[margin=1in]{geometry}
\usepackage[T1]{fontenc}
\usepackage[utf8]{inputenc}
\usepackage{amsmath,amssymb,amsthm}
\usepackage{graphicx}
\usepackage{xcolor}
\usepackage{microtype}
\usepackage[colorlinks=true,linkcolor=blue!55!black,citecolor=blue!55!black,urlcolor=blue!55!black]{hyperref}

\newtheorem{theorem}{Theorem}
\newtheorem{corollary}{Corollary}
\newtheorem{remark}{Remark}
\newcommand{\M}{\mathcal M}
\newcommand{\Lcal}{\mathcal L}

\title{Exact spectral-projection traces of classical unitary propagators}
\author{Candidate full-answer packet for arXiv:2302.00721}
\date{11 August 2026}

\begin{document}
\maketitle

\begin{center}
\fbox{\parbox{0.92\textwidth}{\textbf{Status: candidate full answer to the
trace problem in Remark~3.5, likely valid, subject to human review.}
With the distribution-function convention used in the source,
$U_t=e^{it\Lcal}$ has $|U_t|=I$, so its spectral-projection trace is the
two-step function $\tau(I)\mathbf 1_{[0,1)}$.  Thus the associated weak
Lorentz norm is infinite when $\tau(I)=\infty$ and is independent of $t$ and
$\Lcal$ when $\tau(I)<\infty$.  The alternate literal normal-spectral-measure
reading only counts the resonant eigenspace $e^{it\Lcal}=1$.  Neither reading
can encode classical dispersive decay.}}
\end{center}

\section{Source problem and the notational issue}

G\'omez Cobos, Restrepo, and Ruzhansky study propagators affiliated with the
right group von Neumann algebra $VN_R(G)$ of a locally compact separable
unimodular group.  In Remark~3.5 of arXiv:2302.00721v3 they observe that their
fractional wave and Schr\"odinger estimates do not cover the classical
endpoints and single out the following task \cite{source}:
\begin{quote}
understand the behavior of
$\tau(E_{(\lambda,+\infty)}(e^{it\Lcal}))$ for the classical
Schr\"odinger (wave) propagator; this remains open.
\end{quote}
The source page is reproduced in Appendix~\ref{app:evidence}.

There is a small but decisive ambiguity.  The source defines the distribution
function of an affiliated operator $T$ by
\begin{equation}\label{eq:distribution}
 d_\lambda(T)=\tau\!\left(E_{(\lambda,+\infty)}(|T|)\right),
 \qquad \lambda\ge0,
\end{equation}
where $|T|=(T^*T)^{1/2}$.  This is the projection relevant to generalized
singular numbers and noncommutative Lorentz norms.  Remark~3.5 omits the
modulus, even though $e^{it\Lcal}$ is unitary rather than positive.  We compute
both the intended distribution projection and the literal spectral projection
of the normal operator.  Both are elementary and both obstruct the suggested
route to dispersion.

\section*{Idea of the proof}

Unitary propagation changes phase but not size.  Consequently the modulus of
the propagator is exactly the identity, so every singular value equals one.
This makes its distribution trace completely independent of time and of the
generator.  If one instead takes the spectral projection of the unitary
operator itself, a positive real interval meets the unit circle only at the
point $1$; functional calculus then pulls this point back to the resonant
levels $t\lambda\in2\pi\mathbb Z$ of the generator.  The first interpretation
forgets all phase, while the second sees only exact resonances.  Classical
dispersion comes from cancellation between phases and therefore lies between
these two phase-insensitive extremes.

\section{Exact distribution function and Lorentz norm}

Let $(\M,\tau)$ be a semifinite von Neumann algebra, and let $\Lcal$ be a
positive self-adjoint operator affiliated with $\M$.  The bounded Borel
functional calculus gives $U_t=e^{it\Lcal}\in\M$ for every $t\in\mathbb R$.
Write $T=\tau(I)\in(0,+\infty]$.

\begin{theorem}[complete distribution computation]\label{thm:distribution}
For every $t\in\mathbb R$ and $\lambda\ge0$,
\begin{equation}\label{eq:step}
 d_\lambda(U_t)
 =\tau\!\left(E_{(\lambda,+\infty)}(|U_t|)\right)
 =
 \begin{cases}
 T,&0\le\lambda<1,\\
 0,&\lambda\ge1.
 \end{cases}
\end{equation}
The generalized singular numbers are
\[
 \mu_s(U_t)=
 \begin{cases}
 1,&0<s<T,\\
 0,&s\ge T,
 \end{cases}
\]
with the second line absent when $T=+\infty$.  Hence, for every
$1\le r<\infty$,
\begin{equation}\label{eq:weaknorm}
 \|U_t\|_{L^{r,\infty}(\M)}=T^{1/r},
\end{equation}
where $+\infty^{1/r}=+\infty$.
\end{theorem}

\begin{proof}
The spectral theorem gives $U_t^*=e^{-it\Lcal}$ and therefore
\[
 U_t^*U_t=e^{-it\Lcal}e^{it\Lcal}=I.
\]
Thus $|U_t|=(U_t^*U_t)^{1/2}=I$.  The spectral measure of $I$ is supported at
$1$, so $E_{(\lambda,+\infty)}(I)$ equals $I$ for $0\le\lambda<1$ and equals
$0$ for $\lambda\ge1$.  Taking the trace proves \eqref{eq:step}.

By definition,
$\mu_s(U_t)=\inf\{\lambda\ge0:d_\lambda(U_t)\le s\}$.  Formula
\eqref{eq:step} gives the displayed singular-number formula.  Finally,
\[
 \|U_t\|_{L^{r,\infty}(\M)}
 =\sup_{s>0}s^{1/r}\mu_s(U_t)=T^{1/r},
\]
with the supremum approached from below when $T<\infty$.
\end{proof}

\begin{corollary}[sharp obstruction to the proposed route]\label{cor:route}
If $\tau(I)=+\infty$, then $e^{it\Lcal}\notin L^{r,\infty}(\M)$ for every
finite $r$.  If $\tau(I)<+\infty$, its weak Lorentz norm is the constant
$\tau(I)^{1/r}$, independent of $t$ and $\Lcal$.  The same conclusions hold
for the half-wave unitary $e^{it\sqrt{\Lcal}}$.
\end{corollary}

\begin{proof}
Apply Theorem~\ref{thm:distribution} first to $\Lcal$ and then to the positive
self-adjoint affiliated operator $\sqrt{\Lcal}$.
\end{proof}

\section{The alternate literal spectral-projection reading}

Let $F_t$ denote the projection-valued spectral measure of the normal
operator $U_t$.  Now regard $(\lambda,+\infty)$ as a subset of the real axis
inside $\mathbb C$.

\begin{theorem}[literal normal spectral projection]\label{thm:literal}
Suppose $t\ne0$.  For $0\le\lambda<1$,
\begin{equation}\label{eq:resonance}
 F_t((\lambda,+\infty))
 =F_t(\{1\})
 =E^{\Lcal}\!\left(\{s\ge0:ts\in2\pi\mathbb Z\}\right),
\end{equation}
where $E^{\Lcal}$ is the spectral measure of $\Lcal$.  For $\lambda\ge1$ the
projection is zero.  At $t=0$, it is $I$ below threshold $1$ and zero at or
above $1$.
\end{theorem}

\begin{proof}
The spectrum of $U_t$ lies in the unit circle.  For $0\le\lambda<1$, the real
interval $(\lambda,+\infty)$ meets that circle only at $1$; for
$\lambda\ge1$ the intersection is empty.  Since the spectral measure of
$e^{it\Lcal}$ is the pushforward of $E^{\Lcal}$ under $s\mapsto e^{its}$,
the preimage of $\{1\}$ is exactly the set in \eqref{eq:resonance}.  The case
$t=0$ is immediate from $U_0=I$.
\end{proof}

This literal quantity is well-defined, but it is not the distribution
function \eqref{eq:distribution} used by the source's Lorentz-space method.
It measures only atomic spectral mass at exact resonances and has no general
relation to dispersive cancellation.

\section{Euclidean stress test: infinity or zero, but dispersion survives}

Take $G=\mathbb R^n$ and $\Lcal=-\Delta$.  Under Fourier transform,
$VN_R(\mathbb R^n)$ is the multiplication algebra on frequency space,
$\Lcal$ has symbol $|\xi|^2$, and the Plancherel trace is Lebesgue measure
(up to the harmless Fourier normalization).  Consequently $\tau(I)=+\infty$.
Theorem~\ref{thm:distribution} gives
\[
 \tau\!\left(E_{(\lambda,+\infty)}(|e^{it\Lcal}|)\right)=+\infty
 \quad(0\le\lambda<1),
\]
and the weak $L^r$ multiplier norm is infinite for finite $r$.

For the literal reading and $t\ne0$, the resonant frequency set is
\[
 \{\xi\in\mathbb R^n:t|\xi|^2\in2\pi\mathbb Z\},
\]
a countable union of spheres (and possibly the origin), hence a Lebesgue-null
set.  Its trace is zero.  Thus the two readings produce opposite degenerate
answers: infinity and zero.

Nevertheless, direct Fourier inversion gives the classical Schr\"odinger
kernel, whose absolute value is $(4\pi|t|)^{-n/2}$ for $t\ne0$, and hence
\[
 \|e^{it(-\Delta)}f\|_\infty
 \le (4\pi|t|)^{-n/2}\|f\|_1
\]
with the corresponding choice of Fourier normalization.  This does not
contradict Theorem~\ref{thm:distribution}: the kernel estimate exploits
oscillatory cancellation, whereas generalized singular numbers take the
modulus of the frequency multiplier and discard phase.  It decisively shows
that the proposed trace is not merely difficult to estimate; it is the wrong
invariant for classical dispersion.

\section{Scope, novelty check, and verification}

Theorems~\ref{thm:distribution} and \ref{thm:literal} completely determine
both natural meanings of the trace expression in Remark~3.5.  They answer the
stated trace problem and rule out using that quantity alone to extend the
source's Lorentz-multiplier proof to the classical unitary endpoints.  They do
\emph{not} classify $L^p$--$L^q$ estimates for classical propagators on all
locally compact groups; such estimates require additional structure capable
of retaining phase or kernel cancellation.

The proof was checked directly against Definition~2.4 and Remark~3.5 of the
current arXiv v3 source.  The Euclidean model independently stress-tests both
interpretations.  The run registry and the solution, attempt, and proof-gap
indexes were searched for the arXiv id and central trace/unitary-propagator
phrases.  Bounded exact-title, exact-phrase, arXiv-id, and close-variant web
searches through 11 August 2026 found the source (last revised 16 February
2026, still stating the problem as open) and related spectral-multiplier work,
but no later answer making this computation.  This is not an exhaustive
priority claim.

\textbf{Human-review recommendation.} Verify first that the intended
projection in Remark~3.5 is the distribution projection from Definition~2.4;
then $|e^{it\Lcal}|=I$ makes the main conclusion immediate.  If the omission
of the modulus was deliberate, verify the pushforward identity
\eqref{eq:resonance}.  The principal review judgment is whether to label this
a full answer to the stated trace problem or an obstruction to the broader,
unstated goal of obtaining all classical dispersive estimates.

\appendix
\section{Source-page evidence}\label{app:evidence}

\begin{figure}[p]
\centering
\includegraphics[height=0.82\textheight,keepaspectratio]{figures/open_problem_page.png}
\caption{Page 14 of a local compilation of the official arXiv:2302.00721v3
source.  Remark~3.5 states the classical propagator trace problem.}
\end{figure}

\begin{thebibliography}{1}
\bibitem{source}
Santiago G\'omez Cobos, Joel E. Restrepo, and Michael Ruzhansky,
\emph{Boundedness of spectral multipliers on locally compact groups and
applications}, arXiv:2302.00721v3, 2026.
\end{thebibliography}

\end{document}
